\chapter{\MakeUppercase{Results and Discussion}}

\section{Experimental Results}
\paragraph{}
The implementation of the Secure Messaging App successfully demonstrated the viability of integrating a post-quantum cryptographic architecture within a standard web-desktop hybrid ecosystem. Evaluation of the final build yielded several key results indicating that the system's core requirements were fully met:

\begin{itemize}
    \item \textbf{Successful Hybrid Key Exchange:} The client accurately derived identical master session keys by merging classical NIST P-384 shared secrets with ML-KEM-768 encapsulated outputs via the HKDF algorithm. Integration tests confirmed that any mismatch in the Key Derivation pipeline consistently threw valid errors, preventing unsafe session establishment.
    \item \textbf{Zero-Knowledge Validation:} Traffic interception at the Fastify server node confirmed that payloads transmitted during active sessions consisted entirely of AES-GCM-256 ciphertexts, public initialization vectors (IVs), and authentication tags. The server had zero cryptographic context required to decipher conversational data.
    \item \textbf{Operational UX:} The Electron wrapper successfully maintained strict Context Isolation barriers while delivering a fluid user interface via React. The latency introduced by generating and decapsulating Kyber768 lattice keys was computationally heavier than classical equivalent approaches but remained well within acceptable thresholds (sub-100ms on modern consumer hardware), thereby not noticeably impeding the real-time chat experience.
\end{itemize}

\section{Discussion and Limitations}
\paragraph{}
While the system successfully mitigates the "Harvest Now, Decrypt Later" (HNDL) paradigm and maintains forward secrecy, the integration of post-quantum standards involves inherent technical tradeoffs that warrant discussion.

\begin{itemize}
    \item \textbf{Payload and Storage Overhead:} ML-KEM-768 inherently produces public keys and ciphertexts that are significantly larger than classical ECC keys. The 1,088-byte public keys and 1,088-byte ciphertext encapsulations add non-trivial overhead to the websocket payloads and the Fastify SQL database during the initial prekey polling phase. While modern bandwidth handles this reliably, it is undeniably heavier than traditional Signal Protocol payloads.
    \item \textbf{Implementation Throughput:} Using a pure TypeScript implementation of FIPS 203 ensures maximum cross-platform compatibility across various Electron/OS build targets without relying on complicated native C-bindings. However, this interpreted execution limits raw cryptographic throughput compared to native assembly implementations. Although unnoticeable to individual chat users, this could present challenges for server-side operations if the zero-knowledge constraint were modified to require server-assisted re-encryption for millions of concurrent users.
    \item \textbf{Scalability vs. Security:} The current implementation excels in 1-to-1 secure communication. However, adapting this strict, post-quantum hybrid mechanism for asynchronous group messaging presents a significant scalability challenge (such as extending the Sender Keys protocol to accept lattice dimensions) which was deliberately kept outside the current project boundaries.
\end{itemize}

\section{Conclusion}
\paragraph{}
The results confirm that the Secure Messaging App fundamentally satisfies its primary directive. It proactively neutralizes the quantum threat to digital privacy by deploying a defense-in-depth protocol. By marrying the time-tested discrete logarithm security (ECC) with the newest NIST cryptographic standards (ML-KEM), the system offers robust, future-proof end-to-end encryption accessible via modern, user-friendly desktop frameworks.
